% !TEX root = ../main.tex
%% Devoir surveillé n° 5 — chapitres 14 à 16 (valeurs vérifiées : verifications/devoirs.py)

\devoirDS{5}{Chapitres 14 à 16}{{Le vélo à assistance électrique : puissance et énergie électriques}/7,{La descente à ski : travail et énergie cinétique}/7,{La balle qui rebondit : énergie mécanique, en Python}/6}

\begin{exoBac}{7}{Le vélo à assistance électrique}
Un vélo à assistance électrique est équipé d'une batterie et d'un moteur qui aide le cycliste jusqu'à \qty{25}{km/h}.
\begin{docu}{1}{Caractéristiques}
Batterie : tension nominale \qty{36}{V}, capacité \qty{10}{A.h} ; elle est modélisée par une source réelle de tension, de force
électromotrice $E=\qty{37,0}{V}$ et de résistance interne $r=\qty{0,15}{\ohm}$. Moteur : puissance électrique reçue \qty{250}{W} à pleine
assistance, rendement \qty{80}{\percent}. Prix de l'électricité : 0,25~€ le kilowattheure.
\end{docu}
\begin{questions}
\item Calculer l'énergie stockée dans la batterie chargée, en wattheures puis en joules.
\item À pleine assistance, la batterie fournit \qty{250}{W} sous \qty{36}{V}. Calculer l'intensité du courant.
\item Vérifier, à l'aide du modèle de la batterie, que la tension à ses bornes vaut bien environ \qty{36}{V}. Calculer la puissance
dissipée par effet Joule dans la batterie.
\item Calculer la puissance mécanique fournie par le moteur. Que devient le reste de la puissance reçue ?
\item Calculer l'autonomie à pleine assistance, puis la distance parcourue à \qty{25}{km/h}.
\item Combien coûte une recharge complète de la batterie ?
\end{questions}
\end{exoBac}

\begin{exoBac}{7}{La descente à ski}
Un skieur de masse $m=\qty{70}{kg}$ (équipement compris) part sans vitesse du point A, en haut d'une piste rectiligne, et descend
jusqu'au point B. Sa vitesse mesurée en B vaut $v_\mathrm{B}=\qty{20,0}{m.s^{-1}}$.
\schema{ds5-piste}
\begin{donnees}
$g=\qty{9,81}{N.kg^{-1}}$. On suppose que la force de frottement exercée par la neige et l'air est constante et opposée au mouvement.
\end{donnees}
\begin{questions}
\item Calculer le dénivelé $h$ entre A et B.
\item Calculer le travail du poids du skieur entre A et B.
\item En l'absence de frottement, quelle serait la vitesse du skieur en B ? L'exprimer en \unit{km/h}.
\item Calculer l'énergie cinétique réelle du skieur en B, puis le travail de la force de frottement entre A et B.
\item En déduire la valeur de la force de frottement.
\item Calculer la variation de l'énergie mécanique du skieur et le pourcentage d'énergie perdu. Sous quelle forme cette
énergie est-elle transférée ?
\end{questions}
\end{exoBac}

\begin{exoBac}{6}{La balle qui rebondit}
Une balle de tennis de masse $m=\qty{58}{g}$ est lâchée sans vitesse d'une hauteur de \qty{1,500}{m}. Une vidéo donne l'altitude $y$ de la
balle toutes les \qty{0,05}{s} pendant sa chute ; un programme exploite ces mesures ($g=\qty{9,81}{N.kg^{-1}}$).
\begin{python}
m, g, dt = 0.058, 9.81, 0.05
y = [1.500, 1.488, 1.451, 1.390, 1.304, 1.193, 1.059]     # m
for i in range(1, len(y) - 1):
    v = (y[i - 1] - y[i + 1]) / (2 * dt)
    Ec = 0.5 * m * v**2
    Ep = m * g * y[i]
    print(round(Ec, 3), round(Ep, 3), round(Ec + Ep, 3))
\end{python}
\begin{console}
0.007 0.847 0.854
0.028 0.826 0.853
0.063 0.791 0.854
0.113 0.742 0.854
0.174 0.679 0.853
\end{console}
\begin{questions}
\item \capacite{Capacité numérique}Que calcule la ligne \code{v = ...} ? Pourquoi la boucle ne traite-t-elle ni la première ni la dernière position ?
\item Que représentent les trois nombres affichés sur chaque ligne ? Retrouver par le calcul la valeur 0.847.
\item Comment évolue l'énergie mécanique de la balle pendant la chute ? Conclure.
\item Après le premier rebond, la balle remonte jusqu'à \qty{1,10}{m}. Calculer l'énergie mécanique perdue lors du rebond et la
fraction de l'énergie initiale qu'elle représente.
\item Que devient l'énergie perdue ?
\end{questions}
\end{exoBac}
